% ============================================================================
% Fichier  : partie3/P03_C07_loi_2010.tex
% Titre    : La Loi n° 2010/012 et les Textes Connexes
% Auteur   : MINKA MI NGUIDJOÏ Thierry Emmanuel
% Version  : 3.0 -- novembre 2026
% Corrigés : corriges/Corr_C07.tex
% ============================================================================
% Sources : loi n° 2010/012 du 21 décembre 2010 (texte de la bibliothèque du
% cours), articles 13, 14, 24 à 32, 35, 41 à 51, 74, 91 : cités mot pour mot
% quand ils sont entre guillemets. Loi n° 2024/017, texte officiel. Les
% textes connexes du Tableau 7.4 portent chacun leur degré de vérification ;
% ceux qui sont marqués « à vérifier » sont cités de mémoire et ne doivent
% pas servir de critère d'audit avant lecture du texte.
% ============================================================================

\chapter{La Loi n° 2010/012 et les Textes Connexes}
\label{chap:loi2010}
\chaptermark{La loi n° 2010/012}

\epigraph{%
    « Une loi nouvelle n'efface pas les anciennes.\\
    Elle oblige à les relire. »%
}{Principe du cours}

\niveauChapitre{Fondamental}{%
    Chapitre~\ref{chap:loi}~; droit du numérique (UE MSN-2S3-JUR01)
}

\begin{boxobjectifs}
\begin{itemize}
    \item Présenter ce que la loi n°~2010/012 impose en matière de sécurité,
          de conservation et de vie privée.
    \item Confronter la conservation décennale de la loi de 2010 à la durée
          limitée de la loi de 2024.
    \item Rapprocher les délits de l'article~74 de ceux de la loi de 2024.
    \item Recenser les textes connexes applicables à une organisation, avec
          leur degré de vérification.
    \item Choisir les critères d'un audit et dresser la carte
          d'applicabilité d'une organisation.
\end{itemize}
\end{boxobjectifs}

\bigskip

% ============================================================================
\section*{Avant de commencer : « dix ans, c'est la loi »}
% ============================================================================

Bafoussam, mercredi 5~août~2026. Mme~BELLO demande à M.~KAMGA combien de
temps NYANGA CRÉDIT garde les journaux de connexion de son application
mobile. « Dix ans. C'est la loi de 2010. » Il sort le texte et montre
l'article~29.

Me~ETOUNDI, qui assiste à l'entretien, ouvre l'autre loi. L'article~13 de la
loi de 2024 limite la conservation à « la période nécessaire » au
traitement. Les journaux contiennent l'adresse IP et l'identifiant de chaque
client à chaque connexion. Dix ans de ces données sont-ils nécessaires~?

Les deux ont un texte en main. L'auditrice note que la question n'est pas de
savoir qui a raison, mais de savoir si l'article~29 vise bien un
établissement de microfinance et ce que l'établissement a écrit pour
justifier son choix. Il n'a rien écrit.

\begin{boxalert}
\textbf{Violation typique.} Auditer la protection des données contre la
seule loi n°~2024/017, comme si elle avait remplacé tout ce qui la
précédait.

\smallskip
\textbf{Règle~:} l'article~75 de la loi de 2024 n'abroge que les dispositions
contraires. L'auditeur recense les textes antérieurs encore applicables, les
inscrit parmi ses critères et signale les points où deux textes se
rencontrent sans qu'une règle dise lequel l'emporte.
\end{boxalert}

% ============================================================================
\section{Une Loi de Sécurité qui Parle de Données}
\label{sec:c07-architecture}
% ============================================================================

La loi n°~2010/012 du 21~décembre~2010~\cite{loi2010012} régit la sécurité des
réseaux de communications électroniques et des systèmes d'information, et
réprime la cybercriminalité. Elle n'est pas une loi de protection des
données. Elle en a pourtant tenu lieu pendant quatorze ans, par quelques
articles dispersés.

\begin{tableaucours}{cL{5.4cm}cL{4.0cm}}{%
    \tableauHeader{Titre} &
    \tableauHeader{Objet} &
    \tableauHeader{Art.} &
    \tableauHeader{Intérêt pour ce cours}
}
I & Dispositions générales, définitions & 1--5 & Vocabulaire~: données de
connexion, données de trafic \\
II & Cybersécurité~: régulation, audit, certification électronique,
protection des réseaux, des systèmes et de la vie privée & 6--51 & Sécurité
obligatoire, conservation, vie privée, interception \\
III & Cybercriminalité~: procédure, infractions & 52--89 & Enquête, délits
sur les données (art.~74) \\
IV & Coopération et entraide judiciaire internationales & 90--94 & Demandes
venues de l'étranger \\
V & Dispositions transitoires et finales & 95 et s. & Textes
d'application \\
\end{tableaucours}
\vspace{-0.8em}
\begin{center}{\small\itshape Tableau~7.1. La loi n°~2010/012 en cinq
titres}\end{center}

La loi confie la régulation de la sécurité électronique à l'Agence nationale
des technologies de l'information et de la communication (article~7). Tant
que l'Autorité de protection des données n'est pas organisée, l'Agence est
le seul régulateur en place dont les compétences touchent aux données
personnelles, par le biais de la sécurité.

% ============================================================================
\section{La Sécurité Obligatoire}
\label{sec:c07-securite}
% ============================================================================

Trois obligations de la loi de 2010 s'ajoutent à l'article~27 de la loi de
2024.

L'article~26 demande aux exploitants des systèmes d'information de se doter
« de systèmes normalisés leur permettant d'identifier, d'évaluer, de traiter
et de gérer continûment les risques liés à la sécurité ». C'est une
obligation de gestion des risques, que l'ISO/IEC~27001 permet de satisfaire.

L'article~30 leur demande d'évaluer et de réviser leurs mesures « en fonction
de l'évolution des technologies ».

Les articles~13 et~32 instituent un audit de sécurité obligatoire. Il vise
les réseaux et les systèmes d'information des opérateurs, des autorités de
certification et des fournisseurs de services de communications
électroniques. Il est périodique, conduit par l'Agence ou sous son contrôle
et ses rapports sont confidentiels (article~32).

\begin{boxinfo}
\textbf{Ce que l'auditeur de protection des données en fait.} Il ne refait
pas l'audit de sécurité. Il demande si l'organisation y est assujettie, la
date du dernier audit et la suite donnée à ses recommandations. Le rapport
est confidentiel~: l'auditeur peut ne recevoir qu'une attestation. Il cote
alors l'existence de l'audit, pas son contenu. Le cours d'audit de
cybersécurité (UE MSN-2S3-SEC08) traite l'audit lui-même.
\end{boxinfo}

% ============================================================================
\section{La Conservation Décennale}
\label{sec:c07-conservation}
% ============================================================================

Quatre articles de la loi de 2010 imposent de conserver des données pendant
dix ans.

\begin{tableaucours}{cL{4.2cm}L{6.0cm}}{%
    \tableauHeader{Art.} &
    \tableauHeader{Qui} &
    \tableauHeader{Quoi}
}
25 & Opérateurs de réseaux, fournisseurs de services de communications
électroniques & Données de connexion et de trafic \\
29 & Exploitants des systèmes d'information & Données de connexion et de
trafic de leurs systèmes \\
35 & Fournisseurs d'accès et hébergeurs & Données d'identification des
auteurs de contenus \\
46 & Fournisseurs de contenus & Contenus et données stockées \\
\end{tableaucours}
\vspace{-0.8em}
\begin{center}{\small\itshape Tableau~7.2. Les quatre obligations de
conservation pendant dix ans}\end{center}

\begin{boxjuris}
\textbf{\loicam{n°~2010/012}, article~29, alinéa~1}~\cite{loi2010012}~:
\textit{« Les exploitants des systèmes d'information ont l'obligation de
conserver les données de connexion et de trafic de leurs systèmes
d'information pendant une période de dix (10) ans. »}

\smallskip
\textbf{Portée.} L'alinéa~2 ajoute que ces données « peuvent être
accessibles lors des investigations judiciaires ». La finalité est donc la
preuve pénale. La loi ne définit pas l'exploitant de système d'information~;
on ne sait pas avec certitude si une microfinance ou une clinique en est un
au sens de ce texte.
\end{boxjuris}

Une adresse IP et un identifiant de connexion sont des données personnelles
par la lettre de l'article~5 de la loi de 2024 (Chapitre~\ref{chap:donnees}).
Les journaux visés par l'article~29 sont donc des données personnelles, que
l'article~13 de la loi de 2024 demande de ne pas garder au-delà du
nécessaire.

Les deux textes ne sont pas incompatibles. L'article~9 de la loi de 2024
admet le traitement « nécessaire au respect d'une obligation légale », et
l'article~38 réserve « les exigences de conservation » prévues par d'autres
textes. Une obligation de conserver dix ans, si elle s'applique, fonde la
conservation. La difficulté est en amont~: savoir à qui elle s'applique.

L'auditeur procède en trois temps.

\begin{enumerate}
    \item Il demande à l'organisation si elle se considère comme opérateur,
          fournisseur, hébergeur ou exploitant au sens de la loi de 2010
          et sur quel fondement.
    \item Si elle répond oui, il vérifie que seules les données visées
          (connexion et trafic) sont gardées dix ans, que leur accès est
          restreint et que leur finalité reste la preuve.
    \item Si elle répond non ou ne sait pas, il cote \NEK{} de critère
          l'application de l'article~29 et vérifie que la durée retenue est
          écrite et justifiée au regard de l'article~13.
\end{enumerate}

Dans les trois cas, des journaux gardés dix ans et consultés pour surveiller
les clients ou le personnel sont détournés de leur finalité. Cet écart-là ne
dépend d'aucune interprétation.

% ============================================================================
\section{La Vie Privée dans les Communications}
\label{sec:c07-vieprivee}
% ============================================================================

La section~IV du chapitre~IX (articles~41 à~48) protège la vie privée. Elle
reste applicable en ce qu'elle n'est pas contraire à la loi de 2024 et elle
contient des règles que celle-ci n'a pas reprises.

\begin{tableaucours}{cL{5.6cm}L{4.8cm}}{%
    \tableauHeader{Art.} &
    \tableauHeader{Règle (paraphrasée)} &
    \tableauHeader{Traitement concerné}
}
41 & Droit au respect de la vie privée~; mesures conservatoires du juge &
Tous \\
42 & Confidentialité des communications et des données de trafic &
Messagerie, application, téléphonie \\
44 & Interdiction d'écouter, d'intercepter ou de stocker les communications
sans consentement, sauf autorisation légale & Enregistrement des appels,
supervision de la messagerie \\
45 & Enregistrement autorisé dans le cadre professionnel pour fournir la
preuve d'une communication & Service client, ordres par téléphone \\
46 (2) & Filtres contre les atteintes aux données personnelles & Sites et
applications ouverts au public \\
47 & Consentement préalable pour lire ou écrire dans l'équipement terminal &
Traceurs, kits de mesure d'audience \\
48 & Prospection électronique~: émetteur identifié, adresse valide pour
demander l'arrêt & SMS et courriels commerciaux \\
\end{tableaucours}
\vspace{-0.8em}
\begin{center}{\small\itshape Tableau~7.3. La vie privée dans la loi
n°~2010/012}\end{center}

\begin{boxjuris}
\textbf{\loicam{n°~2010/012}, article~47}~\cite{loi2010012}~:
\textit{« L'utilisation des réseaux de communications électroniques et des
systèmes d'information aux fins de stocker les informations ou d'accéder à
des informations stockées dans un équipement terminal d'une personne
physique ou morale, ne peut se faire qu'avec son consentement préalable. »}

\smallskip
\textbf{Portée.} Un kit de mesure d'audience intégré à une application
mobile, qui lit l'identifiant du téléphone, accède à des informations
stockées dans l'équipement terminal. Il exige un consentement préalable,
distinct de celui que la loi de 2024 demande pour le traitement lui-même.
C'est le cas du traitement que l'audit découvrira dans l'application de
NYANGA (Chapitre~\ref{chap:cartographie}).
\end{boxjuris}

L'article~45 autorise l'enregistrement des communications « effectué dans le
cadre professionnel en vue de fournir la preuve numérique d'une
communication électronique ». L'enregistrement des appels du service client
de NYANGA peut s'appuyer sur ce texte s'il sert à prouver les opérations
convenues par téléphone. Il ne le peut pas s'il sert à évaluer les agents ou
à entraîner un logiciel. La finalité décide, comme à l'article~11 de la loi
de 2024.

% ============================================================================
\section{Interception, Réquisitions et Entraide}
\label{sec:c07-enquete}
% ============================================================================

L'organisation auditée détient des données que les enquêteurs peuvent
demander. La loi de 2010 fixe le cadre de ces demandes. Les officiers de
police judiciaire et les agents assermentés de l'Agence conduisent les
enquêtes selon le Code de procédure pénale (article~52). En cas de crime ou
de délit prévu par la loi, l'officier de police judiciaire peut intercepter
et enregistrer des communications (article~49) et les personnels des
opérateurs sont tenus au secret sur les réquisitions reçues (article~51).
Les demandes venues de l'étranger passent par le ministère chargé des
relations extérieures, sauf convention contraire (article~91)~; la
Convention de Budapest en est une (Chapitre~\ref{chap:international}).

L'auditeur ne juge pas la régularité d'une enquête. Il vérifie que
l'organisation sait recevoir une demande~: qui la reçoit, comment son auteur
et son fondement sont vérifiés, quelles données sont remises, quelle trace
en est gardée. Une remise de données sur simple appel téléphonique d'un
prétendu enquêteur est une violation de données
(Chapitre~\ref{chap:violations}).

% ============================================================================
\section{Deux Lois, des Délits Voisins}
\label{sec:c07-delits}
% ============================================================================

L'article~74 de la loi de 2010 punit plusieurs atteintes aux données
personnelles. Il a été écrit à une époque où aucune formalité préalable
n'existait encore au Cameroun et il en parle pourtant.

\begin{boxjuris}
\textbf{\loicam{n°~2010/012}, article~74, alinéas~3 et~7}~\cite{loi2010012}~:
\textit{« (3) Est puni d'un emprisonnement d'un (01) à trois (03) ans et
d'une amende de 1.000.000 (un million) à 5.000.000 (cinq millions) F CFA ou
de l'une de ces deux peines seulement, quiconque procède ou fait procéder,
même par négligence au traitement des données à caractère personnel en
violation des formalités préalables à leur mise en œuvre. [\ldots]
(7) Est puni d'un emprisonnement de six (06) mois à deux (02) ans et d'une
amende de 5.000.000 (cinq millions) à 50.000.000 (cinquante millions) F
CFA, ou de l'une de ces deux peines seulement, celui qui conserve des
informations sous une forme nominative ou chiffrée au-delà de la durée
légale indiquée dans la demande d'avis ou la déclaration préalable à la mise
en œuvre du traitement automatisé. »}

\smallskip
\textbf{Portée.} L'alinéa~3 punit le traitement fait sans formalité
préalable, « même par négligence ». Depuis 2024, cette formalité existe~:
c'est l'autorisation de l'article~19. L'alinéa~7 punit la conservation
au-delà de la durée déclarée, délit que la loi de 2024 ne reprend pas en ces
termes.
\end{boxjuris}

\begin{tableaucours}{L{4.6cm}cL{4.8cm}c}{%
    \tableauHeader{Comportement} &
    \tableauHeader{2010} &
    \tableauHeader{Loi n°~2024/017} &
    \tableauHeader{2024}
}
Traitement sans formalité préalable & 74 (3) & Traitement sans autorisation
(pénalité administrative) & 55 \\
Collecte par des moyens illicites & 74 (4) & Collecte frauduleuse, déloyale
ou illicite & 63 \\
Données révélant l'origine tribale, les opinions, les mœurs, sans accord
exprès & 74 (5) & Données sensibles sans consentement exprès & 66 \\
Détournement des informations & 74 (6) & Détournement de finalité & 67 \\
Conservation au-delà de la durée déclarée & 74 (7) & Pas de délit
identique~; durée limitée & 13 \\
Divulgation portant atteinte à la considération & 74 (8) & Divulgation
portant atteinte à la dignité ou à la vie privée & 68 \\
Interception de données personnelles en transmission & 74 (2) & Accès
frauduleux & 63 \\
\end{tableaucours}
\vspace{-0.8em}
\begin{center}{\small\itshape Tableau~7.4. Concordance indicative des délits
des deux lois}\end{center}

Les peines diffèrent d'une colonne à l'autre. Quand deux textes punissent le
même fait de peines différentes, lequel s'applique~? L'article~75 de la loi
de 2024 abroge les dispositions « contraires », sans dire si une peine
différente pour le même fait est une contrariété. La réponse appartient au
juge pénal.

L'auditeur n'a pas à la donner. Il rattache l'écart aux deux articles,
signale la concurrence et classe l'incertitude sur la peine parmi les
\NEK{} de critère. L'écart lui-même reste un \NON{}~: les deux lois le
réprouvent.

% ============================================================================
\section{Les Textes Connexes}
\label{sec:c07-connexes}
% ============================================================================

D'autres textes touchent aux données personnelles d'une organisation
camerounaise. Le tableau en donne la liste de travail du cours, avec pour
chacun ce que l'auteur a pu vérifier.

\begin{tableaucours}{L{4.6cm}L{4.4cm}L{3.0cm}}{%
    \tableauHeader{Texte} &
    \tableauHeader{Ce qu'il apporte} &
    \tableauHeader{Vérification}
}
Loi n°~2010/013 sur les communications électroniques, modifiée en 2015 &
Secret des correspondances, obligations des opérateurs envers les abonnés &
Texte en bibliothèque, à lire \\
Décret n°~2015/3759 sur l'identification des abonnés & Collecte de
l'identité de tout abonné & Texte en bibliothèque, à lire \\
Loi-cadre n°~2011/012 sur la protection du consommateur & Information et
loyauté envers le consommateur & Texte en bibliothèque, à lire \\
Code de procédure pénale, articles~102 et~154 & Secret de l'enquête et de
l'information judiciaire & D'après~\cite{minka2026chainepenale} \\
Code pénal & Secret professionnel, violation de correspondance & À
vérifier \\
Loi sur le secret bancaire & Confidentialité des informations des clients &
À vérifier \\
Règlements de la CEMAC et de la COBAC & Conservation et communication des
données des clients, lutte contre le blanchiment, services de paiement & À
identifier \\
Texte sur la protection des enfants en ligne & Obligations envers les
mineurs & À vérifier \\
Code de déontologie médicale & Secret médical & À vérifier \\
\end{tableaucours}
\vspace{-0.8em}
\begin{center}{\small\itshape Tableau~7.5. Les textes connexes et leur degré
de vérification (novembre~2026)}\end{center}

La dernière colonne n'est pas une précaution de style. Un texte « à
vérifier » ne peut pas figurer parmi les critères d'un audit~: l'auditeur le
lit d'abord dans sa version officielle. S'il ne le trouve pas, il l'écrit et
cote \NEK{} de critère les exigences qui en dépendraient. Le
Chapitre~\ref{chap:secteurs} reprend ces textes secteur par secteur.

% ============================================================================
\section{Choisir ses Critères}
\label{sec:c07-criteres}
% ============================================================================

Cinq chapitres ont présenté des textes de nature et de portée différentes.
Avant d'ouvrir une mission, l'auditeur décide lesquels serviront de
critères. La règle tient en quatre lignes.

\begin{boxprocedure}
\textbf{Choisir les critères d'un audit de protection des données.}
\begin{enumerate}
    \item \textbf{L'obligatoire d'abord.} La loi n°~2024/017, la loi
          n°~2010/012 pour ce qui subsiste, les textes sectoriels et
          communautaires applicables à l'organisation. Ils entrent toujours
          dans les critères.
    \item \textbf{L'étranger opposable ensuite.} Le RGPD pour les seuls
          traitements qui remplissent son article~3.
    \item \textbf{Le contractuel.} Les clauses imposées par un client, un
          prestataire ou un investisseur, si le commanditaire veut qu'elles
          soient auditées.
    \item \textbf{Le volontaire déclaré.} Une norme, un cadre ou un repère
          étranger, retenu comme référentiel ou comme critère de
          substitution et nommé comme tel.
\end{enumerate}
\end{boxprocedure}

Le résultat se présente dans une \termecle{carte d'applicabilité}~: un
tableau qui dit, pour chaque texte, son statut à l'égard de l'organisation.

\begin{tableaucours}{L{4.0cm}L{2.6cm}L{2.6cm}L{2.8cm}}{%
    \tableauHeader{Texte} &
    \tableauHeader{NYANGA} &
    \tableauHeader{MBOLO} &
    \tableauHeader{NOUN PAY}
}
Loi n°~2024/017 & Obligatoire & Obligatoire & Obligatoire \\
Loi n°~2010/012, art.~41 à~48 et~74 & Obligatoire & Obligatoire &
Obligatoire \\
Loi n°~2010/012, art.~29 & À qualifier & À qualifier & À qualifier \\
Textes CEMAC et COBAC & Obligatoires & Sans objet & Obligatoires \\
Secret médical & Sans objet & Obligatoire & Sans objet \\
RGPD & Sans objet & Sans objet & Obligatoire pour NOUN DIASPORA~;
contractuel avec l'hébergeur \\
Cadre du NIST & Sans objet & Sans objet & Contractuel (investisseur) \\
ISO/IEC~27701 & Sans objet & Sans objet & Volontaire (projet de
certification) \\
Convention de Malabo & Comparatif & Comparatif & Comparatif \\
\end{tableaucours}
\vspace{-0.8em}
\begin{center}{\small\itshape Tableau~7.6. La carte d'applicabilité des
trois missions}\end{center}

La carte se joint au plan d'audit (Chapitre~\ref{chap:audit}). Elle protège
l'auditeur contre deux reproches~: avoir oublié un texte applicable, ou
avoir reproché à l'organisation de ne pas respecter un texte qui ne la lie
pas.

% ============================================================================
\section{La Loi de 2010 à NYANGA}
\label{sec:c07-cas}
% ============================================================================

\begin{boxscenario}{Mission NYANGA. Les exigences de la loi n° 2010/012, août 2026}
\textbf{Constatations.}
\begin{itemize}
    \item L'application mobile intègre un kit de mesure d'audience qui lit
          l'identifiant du téléphone dès le premier lancement, avant tout
          écran d'information.
    \item Tous les appels au service client sont enregistrés et gardés deux
          ans. Le chef du service en écoute un échantillon chaque mois pour
          noter les agents.
    \item Les SMS de prospection n'indiquent aucun moyen de demander leur
          arrêt.
    \item Des relevés de compte sont envoyés aux clients depuis la
          messagerie personnelle des agents.
    \item Une analyse des risques de sécurité a été refaite après l'incident
          de 2025 et mise à jour en mars~2026.
    \item Les journaux de connexion de l'application sont gardés dix ans.
          Aucune note ne dit pourquoi. L'établissement n'a jamais fait
          l'objet d'un audit de l'Agence et ne sait pas s'il y est
          assujetti.
\end{itemize}

\textbf{Analyse} (confiance élevée sur les faits, moyenne sur
l'assujettissement). Quatre écarts ne dépendent d'aucune interprétation. Les
trois points incertains tiennent tous à la même question~: NYANGA CRÉDIT
est-elle un exploitant de système d'information au sens de la loi de 2010~?
\end{boxscenario}

\begin{boxlecture}{NYANGA CRÉDIT, exigences de la loi n° 2010/012}
\begin{tableaucours}{L{7.6cm}cL{3.4cm}}{%
    \tableauHeader{Exigence} &
    \tableauHeader{Cote} &
    \tableauHeader{Preuve}
}
Le kit de mesure d'audience recueille un consentement préalable (art.~47) &
\NON{} & Analyse de l'application \\
L'enregistrement des appels sert la seule preuve (art.~45) & \NON{} &
Entretien, grille de notation \\
Les messages de prospection permettent de demander l'arrêt (art.~48) &
\NON{} & Échantillon de 20~SMS \\
Les communications avec les clients sont confidentielles (art.~42) & \NON{} &
Messageries personnelles \\
Les risques de sécurité sont gérés de façon continue (art.~26) & \OUI{} &
Analyse de mars~2026 \\
L'audit de sécurité obligatoire a été réalisé (art.~32) & \NEK{} &
Assujettissement non établi \\
L'obligation de conserver dix ans s'applique (art.~29) & \NEK{} & Qualité
d'exploitant non établie \\
La durée de dix ans est justifiée au regard de l'article~13 de la loi de
2024 & \NEK{} & Dépend de la ligne précédente \\
\end{tableaucours}

\medskip
\textbf{Décompte.} Un \OUI{}, quatre \NON{}, trois \NEK{} sur huit.
\hfill\lecture{0{,}125}{0{,}375}{1{,}41}

\medskip
\textbf{Lecture.} Domaine opaque. Les trois \NEK{} sont de critère et se
lèveraient ensemble si la notion d'exploitant était précisée. Les quatre
\NON{} appellent des mesures immédiates, dont le retrait du kit de mesure
d'audience.
\end{boxlecture}

% ============================================================================
\section{Ce que les Textes demandent à chaque Métier}
\label{sec:c07-metiers}
% ============================================================================

\begin{boxmetier}{Juriste conformité}
Tenir la liste des textes applicables avec leur degré de vérification.
Écrire, pour chaque point où deux textes se rencontrent, la lecture retenue
par l'organisation et sa justification.
\end{boxmetier}

\begin{boxmetier}{RSSI}
Séparer les journaux gardés pour la preuve de ceux qui servent
l'exploitation. Restreindre l'accès aux premiers et documenter leur durée.
\end{boxmetier}

\begin{boxmetier}{Auditeur de la protection des données}
Joindre la carte d'applicabilité au plan d'audit. Rattacher chaque écart à
tous les articles qui le visent, dans les deux lois.
\end{boxmetier}

\begin{boxresume}
\begin{itemize}
    \item La loi n°~2010/012 reste applicable pour tout ce qui n'est pas
          contraire à la loi de 2024~: sécurité, conservation, vie privée
          dans les communications, délits.
    \item Elle impose une gestion continue des risques de sécurité
          (article~26) et un audit de sécurité obligatoire pour certains
          acteurs (articles~13 et~32).
    \item Quatre articles imposent de conserver des données dix ans. Ils
          fondent la conservation s'ils s'appliquent~; il faut d'abord
          établir à qui.
    \item L'article~47 exige un consentement préalable pour tout accès à
          l'équipement terminal~; l'article~45 limite l'enregistrement des
          communications à la preuve.
    \item L'auditeur choisit ses critères dans l'ordre~: obligatoire,
          étranger opposable, contractuel, volontaire déclaré. Il les
          présente dans une carte d'applicabilité.
\end{itemize}
\end{boxresume}

\begin{boxinfo}
\textbf{Points de contrôle.} Vous avez compris ce chapitre si vous savez~:
\begin{itemize}
    \item citer trois règles de la loi de 2010 absentes de la loi de 2024~;
    \item expliquer comment coter la conservation de journaux pendant dix
          ans~;
    \item dresser la carte d'applicabilité d'une organisation.
\end{itemize}
\end{boxinfo}

\begin{boxexo}[title={\ding{45}\quad Exercice~7.1. Quel article \niveaubase}]
Indiquez l'article de la loi n°~2010/012 qui s'applique~: (a) une
application installe un traceur sans rien demander~; (b) un établissement
envoie des SMS publicitaires sans numéro pour se désinscrire~; (c) un
responsable écoute les appels de ses agents à leur insu~; (d) une entreprise
garde les enregistrements d'appels comme preuve des ordres reçus.
\end{boxexo}

\begin{boxexo}[title={\ding{45}\quad Exercice~7.2. Les appels enregistrés \niveaumoyen}]
Le service client de NYANGA enregistre tous les appels, les conserve deux
ans et les écoute pour évaluer les agents. Analysez ce traitement au regard
des articles~44 et~45 de la loi n°~2010/012 et des articles~9 et~11 de la
loi n°~2024/017. Proposez une règle conforme.
\end{boxexo}

\begin{boxexo}[title={\ding{45}\quad Exercice~7.3. La carte de MBOLO \niveaumoyen}]
Le CENTRE MÉDICAL MBOLO envisage d'accueillir des patients envoyés par une
organisation humanitaire établie à Genève, qui lui impose par contrat ses
propres règles de protection des données. Refaites la colonne MBOLO du
Tableau~7.6. Quels textes changent de statut et lesquels restent à
vérifier avant d'entrer dans les critères~?
\end{boxexo}

\begin{boxexo}[title={\ding{45}\quad Exercice~7.4. Dix ans ou le nécessaire \niveauexpert}]
M.~KAMGA vous demande une réponse écrite~: NYANGA doit-elle garder ses
journaux de connexion dix ans~?
\begin{enumerate}[(a)]
    \item Rédigez la note, en séparant ce qui est établi par les textes de
          ce qui relève d'une interprétation.
    \item Proposez une politique de conservation des journaux qui soit
          défendable dans les deux lectures de l'article~29.
    \item Dites quelles cotes de l'encadré de lecture changeraient si
          l'Agence confirmait par écrit que l'établissement est un
          exploitant au sens de la loi.
\end{enumerate}
\end{boxexo}

\bigskip

\begin{boxinfo}
\textbf{Pour aller plus loin.}
\begin{itemize}
    \item La loi n°~2010/012~\cite{loi2010012}, articles~24 à~51 et~74, dans
          son texte officiel.
    \item T.~E.~Minka Mi Nguidjoï, R.~T.~Massoussi Songo et
          D.~Bialo~\cite{minka2026chainepenale}, pour l'articulation des
          deux lois dans la chaîne pénale.
\end{itemize}
\end{boxinfo}

\begin{boxinfo}
\textbf{Transition.} Le contexte est posé~: ses sources, le règlement
européen, les référentiels volontaires, les deux lois camerounaises et la
manière de choisir entre eux. Il reste à s'en servir. La Partie~IV ouvre la
mission et commence par la manière de conduire un audit. C'est l'objet du
Chapitre~\ref{chap:audit}.
\end{boxinfo}
